\documentclass[12pt,a4paper]{article}
\usepackage[utf8]{inputenc}
\usepackage[T1]{fontenc}
\usepackage{amsmath,amssymb}
\usepackage{enumitem}
\usepackage{geometry}
\usepackage{array}
\usepackage{multicol}
\geometry{margin=2.2cm}
\setlist{itemsep=0.25em, topsep=0.3em}
\title{Aula 09 -- Dependência Funcional e Normalização}
\author{Banco de Dados -- Guia de Revisão}
\date{}
\begin{document}
\maketitle
\tableofcontents
\newpage

\section{Objetivo da aula}
Esta aula explica por que alguns esquemas relacionais são ruins mesmo quando funcionam. O foco é reconhecer redundância, anomalias, dependências funcionais e formas normais. Para a prova, o mais importante é saber analisar uma relação dada e responder: quais dependências existem, qual é a chave, qual forma normal é violada e qual decomposição corrige o problema.

\section{Ideia central}
Normalização é o processo de organizar atributos em relações para reduzir redundância e evitar anomalias. A decisão não é estética: ela depende das dependências funcionais do domínio.

\subsection{Boa semântica de relação}
Uma relação deve agrupar atributos que descrevem uma mesma entidade ou relacionamento. Quando uma tabela mistura conceitos diferentes, aparecem muitos valores repetidos, valores nulos e dependências indiretas.

\subsection{Anomalias}
\begin{itemize}
    \item \textbf{Atualização}: a mesma informação aparece repetida e precisa ser alterada em várias tuplas.
    \item \textbf{Inserção}: não é possível inserir uma informação sem fornecer outra informação ainda inexistente.
    \item \textbf{Remoção}: ao remover uma tupla, perde-se uma informação que deveria permanecer.
\end{itemize}

\section{Dependência funcional}
A dependência funcional $X \rightarrow Y$ significa: se duas tuplas têm o mesmo valor em $X$, então devem ter o mesmo valor em $Y$.

Exemplo:
\[
\text{idCurso} \rightarrow \text{nomeCurso}
\]
Se o id do curso é 10, o nome do curso deve ser sempre o mesmo.

\section{Exemplo visual: redundância}
\[
\begin{array}{|c|c|c|c|}
\hline
\text{idAluno} & \text{nomeAluno} & \text{idCurso} & \text{nomeCurso}\\
\hline
1 & \text{Ana} & 10 & \text{BTI}\\
2 & \text{Bruno} & 10 & \text{BTI}\\
3 & \text{Carla} & 20 & \text{Computacao}\\
\hline
\end{array}
\]

Aqui há redundância porque $\text{idCurso} \rightarrow \text{nomeCurso}$. O nome do curso BTI é repetido para todos os alunos do curso 10.

Decomposição adequada:
\[
Aluno(\text{idAluno}, \text{nomeAluno}, \text{idCurso})
\]
\[
Curso(\text{idCurso}, \text{nomeCurso})
\]

\section{Fechamento de atributos}
O fechamento $X^+$ é o conjunto de atributos que podem ser determinados por $X$ usando as dependências funcionais.

Procedimento:
\begin{enumerate}
    \item Comece com $X^+ = X$.
    \item Se existe $A \rightarrow B$ e $A \subseteq X^+$, adicione $B$ a $X^+$.
    \item Repita até não conseguir adicionar mais nada.
\end{enumerate}

Exemplo:
\[
R(A,B,C,D),\quad F=\{A\rightarrow B,\ B\rightarrow C,\ A\rightarrow D\}
\]
Então:
\[
A^+ = \{A,B,C,D\}
\]
Logo, $A$ é chave se determina todos os atributos da relação.

\section{Chaves candidatas}
Uma chave candidata é um conjunto mínimo de atributos que determina todos os atributos da relação.

Exemplo:
\[
R(A,B,C,D),\quad F=\{AB\rightarrow C,\ C\rightarrow D\}
\]
\[
(AB)^+ = \{A,B,C,D\}
\]
Então $AB$ é chave candidata. $A$ sozinho não determina tudo, e $B$ sozinho também não.

\section{Formas normais}
\subsection{1FN}
A relação está na Primeira Forma Normal quando todos os atributos possuem valores atômicos. Uma célula não deve guardar lista de valores.

Violação:
\[
Aluno(\text{matricula}, \text{nome}, \text{telefones})
\]
se \texttt{telefones} guarda vários telefones na mesma célula.

\subsection{2FN}
A 2FN elimina dependências parciais de atributos não-chave em relação a uma chave composta.

Exemplo de violação:
\[
Matricula(\text{idAluno},\text{idDisciplina},\text{nomeAluno},\text{nomeDisciplina},\text{nota})
\]
Chave: $(\text{idAluno},\text{idDisciplina})$.

Dependências:
\[
\text{idAluno} \rightarrow \text{nomeAluno}
\]
\[
\text{idDisciplina} \rightarrow \text{nomeDisciplina}
\]
Como atributos não-chave dependem só de parte da chave composta, há violação de 2FN.

\subsection{3FN}
A 3FN elimina dependências transitivas de atributos não-chave.

Exemplo:
\[
Funcionario(\text{cpf},\text{nome},\text{idDepto},\text{nomeDepto})
\]
\[
\text{cpf} \rightarrow \text{idDepto},\quad \text{idDepto} \rightarrow \text{nomeDepto}
\]
Logo, $\text{cpf} \rightarrow \text{nomeDepto}$ indiretamente. Isso é dependência transitiva.

\subsection{BCNF}
Uma relação está em BCNF quando, para toda dependência funcional não trivial $X \rightarrow Y$, $X$ é superchave.

Exemplo de violação:
\[
R(A,B,C),\quad A\rightarrow B,\quad B\rightarrow C
\]
Se $A$ é chave, mas $B$ não é superchave, então $B\rightarrow C$ viola BCNF.

\section{Cobertura mínima}
Uma cobertura mínima é uma versão simplificada de um conjunto de dependências funcionais, sem atributos ou dependências redundantes.

Passos típicos:
\begin{enumerate}
    \item Separar lado direito: $A\rightarrow BC$ vira $A\rightarrow B$ e $A\rightarrow C$.
    \item Remover atributos desnecessários do lado esquerdo.
    \item Remover dependências que podem ser inferidas por outras.
\end{enumerate}

\section{Junção sem perdas e preservação}
\textbf{Junção sem perdas}: ao decompor e depois juntar, não aparecem tuplas espúrias.

\textbf{Preservação de dependências}: as dependências originais ainda podem ser verificadas nas relações decompostas sem precisar juntar tudo.

Exemplo:
\[
R(A,B,C),\quad F=\{A\rightarrow B,\ B\rightarrow C\}
\]
Decomposição:
\[
R_1(A,B),\quad R_2(B,C)
\]
Preserva diretamente $A\rightarrow B$ em $R_1$ e $B\rightarrow C$ em $R_2$.

\section{Pegadinhas comuns}
\begin{itemize}
    \item Chave simples não viola 2FN por dependência parcial.
    \item Toda relação em BCNF está em 3FN, mas nem toda 3FN está em BCNF.
    \item Dependência transitiva costuma ser cobrada com três atributos: $A\rightarrow B$ e $B\rightarrow C$.
    \item Redundância nem sempre é erro de SQL; muitas vezes é erro de projeto relacional.
\end{itemize}

\section{Desafios}
\subsection*{Desafio 1}
Considere:
\[
R(A,B,C,D),\quad F=\{A\rightarrow B,\ B\rightarrow C,\ CD\rightarrow A\}
\]
Qual é o fechamento de $CD$? $CD$ é chave?

\subsection*{Desafio 2}
Considere:
\[
Pedido(\text{idPedido},\text{idCliente},\text{nomeCliente},\text{data})
\]
\[
\text{idPedido}\rightarrow \text{idCliente},\text{data}
\]
\[
\text{idCliente}\rightarrow \text{nomeCliente}
\]
Qual forma normal é violada e qual decomposição é adequada?

\section{Resolução dos desafios}
\textbf{Desafio 1}: $CD^+ = \{C,D,A,B\}$. Começa com $C,D$, usa $CD\rightarrow A$, depois $A\rightarrow B$, depois $B\rightarrow C$. Como determina todos os atributos, $CD$ é chave.

\textbf{Desafio 2}: há dependência transitiva: $idPedido\rightarrow idCliente\rightarrow nomeCliente$. Viola 3FN. Decomposição: $Pedido(idPedido,idCliente,data)$ e $Cliente(idCliente,nomeCliente)$.

\section{Checklist final}
Antes da prova, saiba responder:
\begin{itemize}
    \item O que significa $X\rightarrow Y$?
    \item Como calcular fechamento?
    \item Como identificar chave candidata?
    \item Qual a diferença entre 1FN, 2FN, 3FN e BCNF?
    \item Como reconhecer dependência parcial e transitiva?
    \item Como decompor uma relação ruim?
\end{itemize}

\end{document}
